% Solution to AMM Problem 12592 (proposed by Elliot Glazer), for submission.
% Parts (a),(b) follow THM-2225 (dyadic block compression; cyclic checksum).
% The "Note to the editor" section is the part (c) extension: THM-2160 S5
% (half-tail rule, max(n+1,2n-2)) + THM-2253 (dyadic contrast, n+2^{nu_2(n)})
% + the linear-deadline (Cn+D) discussion (HYP-9061 frontier, stated as open).
% Everything below was machine-verified exactly; see
% 04-computation/{dyadic_coin_critical_run_checksum_thm2225,
%   biased_coin_dyadic_half_tail_referee_codex_20260724,
%   categorical_coin_online_dyadic_contrast_thm2253}.py
\documentclass[11pt]{article}
\usepackage[margin=1.1in]{geometry}
\usepackage{amsmath,amssymb,amsthm}
\usepackage{booktabs}

\newtheorem{theorem}{Theorem}
\newtheorem{lemma}[theorem]{Lemma}
\newtheorem{proposition}[theorem]{Proposition}
\theoremstyle{remark}
\newtheorem*{remark}{Remark}

\newcommand{\head}{\textnormal{\textsc{heads}}}
\newcommand{\tail}{\textnormal{\textsc{tails}}}

\title{Solution to Problem 12592}
\author{[SOLVER NAME(S)], [City, State/Country]}
\date{[Month Year]}

\begin{document}
\maketitle

\noindent\textbf{Problem 12592} (proposed by Elliot Glazer)\textbf{.}
\emph{Consider a biased coin $C$ with sides labeled $0$ and $1$. When $C$ is
flipped, it shows $0$ with unknown probability $p$ and $1$ with probability
$1-p$, where $0<p<1$. Let the \emph{critical value} of a nonconstant bit
sequence be the number of bits in the maximal constant string starting at the
beginning.}

\emph{(a) Show how to use $C$ to simulate a fair coin by a method that assigns
to any bit sequence with critical value $n$ a result after at most $2n$ flips
of $C$. For example, if the first five flips of $C$ are $00001$, then $n=4$,
and the method must declare heads or tails after at most three more flips.}

\emph{(b) Find a simulation method as in part (a) that always gives a result
after $\max(2,\,2n-1)$ flips of $C$, where $n$ is the critical value.}

\section*{Notation and the exactness lemma}

Write the flips as independent bits $X_1,X_2,\dots$ with $\Pr(X_i=0)=p$ and
$\Pr(X_i=1)=q:=1-p$. For a nonconstant sequence, the critical value is
$n=\min\{k\ge 1: X_{k+1}\ne X_1\}$. A \emph{simulation method} is an online
rule: after each flip it either stops and declares \head{} or \tail{}, or
flips again, and the decision taken at time $k$ depends only on
$X_1,\dots,X_k$. The method is \emph{exactly fair} if
$\Pr(\head)=\Pr(\tail)=\tfrac12$ for every $p\in(0,1)$.

Both solutions below, and the further construction in the closing note, fit
one frame. For a sequence of critical value $n$, let $m$ be the unique power
of two with
\begin{equation}\label{eq:shell}
m\le n\le 2m-1 ,
\end{equation}
so that the first $m$ flips are constant while the first $2m$ are not. The
word $w=X_1\cdots X_{2m}$ then lies in the \emph{shell}
\[
S_m=\{w\in\{0,1\}^{2m}:\ w_1=\dots=w_m,\ w \text{ nonconstant}\},
\]
and distinct $m$ give disjoint cylinder events that exhaust the nonconstant
sequences.

\begin{lemma}[Exactness]\label{lem:exact}
Suppose a method decides every sequence of critical value $n$ by flip $2m$,
with $m$ as in \eqref{eq:shell}, so that on each shell the outcome is a
function $F_m\colon S_m\to\{\head,\tail\}$. If for every $m$ and every $k$
the words of $S_m$ containing exactly $k$ ones split equally between \head{}
and \tail{}, then the method is exactly fair.
\end{lemma}

\begin{proof}
All words of $S_m$ with $k$ ones have the same probability
$p^{\,2m-k}q^{\,k}$, so the equal split makes the \head{} and \tail{}
probabilities agree shell by shell. The shells are disjoint events whose
union is the event that the sequence is nonconstant, and the two constant
sequences have probability $\lim_N p^N=\lim_N q^N=0$. Summing over shells,
$\Pr(\head)=\Pr(\tail)$, and the two add to $1$.
\end{proof}

Note that a method may reach its verdict before flip $2m$; the lemma applies
as long as the verdict is determined by the first $2m$ flips.

\section*{Solution to (a): dyadic block compression}

\emph{The rule.} Flip until the first disagreement $X_{n+1}\ne X_1$ reveals
$n$, let $m$ be as in \eqref{eq:shell}, and flip on through flip
$2m\ (\le 2n)$. Now compress the word $w=X_1\cdots X_{2m}$: group consecutive
symbols into pairs $(w_1w_2)(w_3w_4)\cdots$. If some pair is unequal, the
\emph{leftmost} unequal pair decides: declare \head{} if it is $01$ and
\tail{} if it is $10$. If all pairs are equal, replace each pair by its
repeated symbol (halving the length) and repeat. A nonconstant word of
power-of-two length can never be compressed to a single symbol, so after some
number $r\ge 0$ of halvings a deciding pair appears. The deciding pair at
depth $r$ is a pair of adjacent constant blocks of $w$ of length $2^r$,
jointly occupying an \emph{aligned} interval of length $2^{r+1}$ (positions
$(k-1)2^{r+1}+1$ through $k2^{r+1}$ for some $k$); the compression reaches
depth $r$ precisely because every aligned block of length $2^r$ is constant
and, at depth $r$, every pair to the left of the deciding one is equal.

\emph{Fairness.} Let $\sigma(w)$ be the word obtained from $w$ by swapping
the two constant blocks of its deciding pair. We claim $\sigma$ is a
fixed-point-free involution of $S_m$ that preserves the number of ones and
reverses the verdict; by Lemma~\ref{lem:exact} the method is then exactly
fair, since $\sigma$ matches the \head{} and \tail{} words of each
composition class in pairs.

The deciding pair cannot lie inside the constant first half, so its aligned
interval, having power-of-two length, either lies inside the second half or
is the whole word; in the latter case $w=b^m\bar b^m$ and
$\sigma(w)=\bar b^m b^m$. In both cases the first half of $\sigma(w)$ is
constant, and $\sigma(w)$ is nonconstant, so $\sigma(w)\in S_m$. Swapping two
constant blocks keeps every aligned block of length $2^r$ constant, so the
compression of $\sigma(w)$ again reaches depth $r$; the pairs strictly to the
left of the deciding interval are untouched and still equal, and the deciding
pair itself is still unequal but reversed. Hence the deciding pair of
$\sigma(w)$ is the same interval with the opposite orientation: the verdict
flips. Finally $\sigma$ permutes the letters of $w$, so it preserves
composition, and $\sigma\circ\sigma=\mathrm{id}$ with no fixed points.

\emph{Deadline.} The verdict comes at flip $2m\le 2n$ by \eqref{eq:shell}.
For the sample $00001$ we have $n=4$, $m=4$: at most three further flips.
\qed

\section*{Solution to (b): the cyclic checksum}

\emph{The rule.} If the first two flips disagree, stop at flip $2$: declare
\head{} on $01$ and \tail{} on $10$. Otherwise $m\ge 2$ in
\eqref{eq:shell}. Write the second half of the shell word as
$y_i=X_{m+i}$ for $1\le i\le m$ and set
\[
s(y)\;=\;\sum_{i=1}^{m} i\,y_i \bmod m,\qquad s(y)\in\{0,1,\dots,m-1\}.
\]
Declare \head{} if $s(y)<m/2$ and \tail{} if $s(y)\ge m/2$.

\emph{Deadline.} The coefficient of $y_m$ in $s$ is $m\equiv 0\pmod m$, so
$s$ is determined by flips $m+1,\dots,2m-1$ together with the value of $m$
--- and $m$ is known as soon as the first disagreement appears. If
$n\le 2m-2$, the disagreement occurs by flip $2m-1$, so the method stops at
flip $2m-1\le 2n-1$ (using $m\le n$). If $n=2m-1$, the first $2m-1$ flips
are constant and flip $2m$ itself is needed, but then
$2m=n+1\le 2n-1$ because $n\ge 3$. With the base case this gives
termination within $\max(2,\,2n-1)$ flips. (For $00001$: $n=4$, $m=4$, and
$s=y_1+2y_2+3y_3 \bmod 4$ with $y_1=1$ already seen, so at most two further
flips are used, one fewer than part (a) requires.)

\emph{Fairness.} By Lemma~\ref{lem:exact} (the verdict is a function of the
first $2m$ flips) it suffices to split every composition class of every shell
evenly. Fix the shell of $m\ge 2$ and a class with $k$ ones, $1\le k\le
2m-1$.

If $k\ne m$, the constant first half is forced ($0^m$ if $k<m$, $1^m$ if
$k>m$), so the class is a copy of the set of tails $y$ of fixed weight
$j\in\{1,\dots,m-1\}$, $j\equiv k\pmod m$. Write $j=2^au$ with $u$ odd and
put $\delta=m/2^{a+1}$, a positive integer since $2^a\le j<m$. Let $\rho$
rotate the tail cyclically, $(\rho y)_{i+1}=y_i$ (indices mod $m$). Then
\[
s(\rho y)\;=\;\sum_{i=1}^{m} (i+1)\,y_i - m\,y_m \;\equiv\; s(y)+j
\pmod m ,
\]
so $s(\rho^{\delta}y)\equiv s(y)+\delta j = s(y)+u\cdot\tfrac m2\equiv
s(y)+\tfrac m2 \pmod m$, as $u$ is odd. Therefore the map sending $y$ to
$\rho^{\delta}y$ when $s(y)<m/2$ and to $\rho^{-\delta}y$ when
$s(y)\ge m/2$ is a fixed-point-free, weight-preserving involution of the
class that exchanges \head{} and \tail{}: the class splits evenly.

If $k=m$, the class consists of exactly two words, $0^m1^m$ and $1^m0^m$,
with checksums
\[
s(1^m)=\tfrac{m(m+1)}2\equiv \tfrac m2, \qquad s(0^m)=0 \pmod m ,
\]
so they receive opposite verdicts. (Stopping at flip $2m-1$ when possible is
harmless: both one-flip extensions of such a stopping word carry the same
verdict, and $p+q=1$.) By Lemma~\ref{lem:exact} the method is exactly fair.
\qed

\bigskip
\begin{center}\rule{0.35\textwidth}{0.4pt}\end{center}

%======================================================================
% The section below is intended for the "Note to the editor" box.
%======================================================================
\section*{Note to the editor: a part (c), and the linear-deadline frontier}

The mechanisms above can be pushed one step further. We record here a
sharper envelope with full proof, the matching small-$n$ optimality, and
what we can and cannot say about the proposer's linear-deadline question
(minimize $C$ such that, for some constant $D$, a method always terminates
within $Cn+D$ flips). Write $T(n)$ for the worst-case number of flips a
given method uses over all sequences of critical value $n$.

\begin{lemma}[Floor]\label{lem:floor}
For every exactly fair method, no stopping word is constant; consequently
$T(n)\ge n+1$ for every $n$. In particular $T(1)\ge2$, $T(2)\ge3$,
$T(3)\ge4$.
\end{lemma}

\begin{proof}
If the method could stop at the constant word $b^k$ with verdict $v$, then
$\Pr(v)\ge \Pr(X_1=\dots=X_k=b)=p^k$ or $q^k$, which exceeds $\tfrac12$ for
$p$ near $0$ or $1$, contradicting exact fairness. Since the first $n$ flips
of a critical-value-$n$ sequence are constant, no verdict can be reached
before flip $n+1$.
\end{proof}

So the envelope $\max(2,2n-1)$ of part (b) is already optimal at $n=1$ and
$n=2$, and \emph{no} method achieves the uniform shift $\max(2,2n-2)$: that
would require deciding critical value $2$ by flip $2$, i.e.\ stopping at the
constant words $00$ and $11$. The correct sharpening carries an exception at
$n=2$ and is best possible at $n=1,2,3$:

\begin{theorem}[a part (c)]\label{thm:c}
There is an exactly fair method with
\[
T(n)\;\le\;\max(n+1,\;2n-2)\qquad\text{for every } n\ge 1,
\]
that is, $T(1)=2$, $T(2)=3$, and $T(n)\le 2n-2$ for all $n\ge3$. Moreover
this method stops within $\tfrac32 n$ flips when $n$ is a power of two,
$n\ge4$, and within $n+1$ flips when $n=2m-1$ in \eqref{eq:shell}. By
Lemma~\ref{lem:floor}, its values at $n=1,2,3$ are exactly optimal.
\end{theorem}

\begin{proof}
For $n=1$ use flips one and two as in part (b). Fix a shell
\eqref{eq:shell} with $m\ge2$ and put $t=m/2$, $b=X_1$, and
$z_i=X_{m+i}\oplus b$ ($1\le i\le m$), the tail relative to the run bit;
membership in the shell says exactly that $z\ne0$, and $z_1=1$ if and only
if $n=m$. The rule has two strata.

\emph{(i) If $z_1=1$} (so $n=m$), stop at flip $m+t=\tfrac32 m$ and declare
\head{} iff $b+z_2+\dots+z_t\equiv 0 \pmod 2$. (For $m=2$ the sum is empty:
$001{*}\mapsto\head$, $110{*}\mapsto\tail$ at flip $3$.)

\emph{(ii) If $z_1=0$} (so $n\ge m+1$), read through flip $2m$ and declare
by a fixed table, constructed as follows. Work on the branch $b=0$ and let
$1\le j\le m-1$. Stratum (i) has already assigned verdicts to the weight-$j$
relative words with $z_1=1$; let $A_j$ of them be \head. Among the
$\binom{m-1}{j}$ weight-$j$ words with $z_1=0$, declare the first
\[
R_j:=\tfrac12\tbinom mj-A_j
\]
of them (in lexicographic order) \head{} and the rest \tail. On the branch
$b=1$, give the relative word $z$ the opposite verdict to the one it
receives on the branch $b=0$.

\emph{The table is legal, i.e.\ $R_j$ is an integer with $0\le R_j\le
\binom{m-1}j$.} Counting stratum-(i) \head{} words of weight $j$ (one for
$z_1$; an even number among $z_2,\dots,z_t$; anything among
$z_{t+1},\dots,z_m$) gives the generating function
\[
\sum_j A_j u^j
= u\,\frac{(1+u)^{t-1}+(1-u)^{t-1}}2\,(1+u)^{t}
= \tfrac12\Bigl[u(1+u)^{m-1}+u(1+u)(1-u^2)^{t-1}\Bigr].
\]
Writing $e_j:=[u^j]\,u(1+u)(1-u^2)^{t-1}$, so that
$e_{2a+1}=e_{2a+2}=(-1)^a\binom{t-1}a$ for $0\le a\le t-1$, we get
$A_j=\tfrac12\bigl[\binom{m-1}{j-1}+e_j\bigr]$ and, by Pascal's identity,
\[
R_j=\tfrac12\Bigl[\tbinom{m-1}{j}-e_j\Bigr].
\]
Both $m-1=2^r-1$ and $t-1=2^{r-1}-1$ have all binary digits equal to one,
so by Lucas's theorem $\binom{m-1}j$ and $\binom{t-1}a$ are odd and $R_j$ is
an integer. For the range, $|e_j|=\binom{t-1}a\le\binom{2t-1}{j}$ for
$j\in\{2a+1,2a+2\}$, $j\le m-1$: by Vandermonde,
$\binom{2t-1}{2a+1}\ge\binom{t-1}a\binom{t}{a+1}$ and
$\binom{2t-1}{2a+2}\ge\binom{t-1}a\binom{t}{a+2}$, and the second factors
are $\ge1$ in the stated ranges (for $j=2a+2\le m-1$ one has $a\le t-2$).
Hence $0\le R_j\le\binom{m-1}j$.

\emph{Fairness.} On branch $b=0$, the weight-$j$ layer has exactly
$A_j+R_j=\tfrac12\binom mj$ \head{} words for $1\le j\le m-1$; branch $b=1$
carries the complementary verdicts, so its layers are evenly split as well.
A composition class of the shell with $k$ ones is: for $1\le k\le m-1$, the
branch-$0$ layer $j=k$; for $m+1\le k\le 2m-1$, the branch-$1$ layer
$j=2m-k$; and for $k=m$, the pair $\{0^m1^m,\,1^m0^m\}$, which consists of
the \emph{same} relative word $z=1^m$ on the two branches and therefore
receives opposite verdicts. Every class splits evenly, and
Lemma~\ref{lem:exact} applies.

\emph{Deadlines.} Stratum (i) uses $m+t=\tfrac32m$ flips: this is $3$ when
$m=2$ ($=n+1$) and at most $2m-2=2n-2$ when $m\ge4$. Stratum (ii) uses $2m$
flips with $n\ge m+1$, so $2m\le 2n-2$; when $n=2m-1$ this reads
$2m=n+1$. Together with $T(1)=2$ this is the claimed envelope.
\end{proof}

For the sample $00001$ ($n=m=4$): the method stops after \emph{one} further
flip, declaring \head{} iff $X_6=0$.

We verified all three constructions by machine, exhaustively through the
shell $2m=32$: exact composition balance in every shell (hence exact
fairness for every $p$), and the stated deadlines on every path.

\medskip
\noindent\textbf{The linear-deadline frame.}
The proposer's commentary proposes minimizing $C$ such that some method has
$T(n)\le Cn+D$ for a constant $D$. Three remarks, in increasing order of
openness.

\emph{1. The literal frame is degenerate at small $n$.}
By Lemma~\ref{lem:floor}, $T(1)\ge2$ forces $C+D\ge2$; so once $C=2$, the
smallest admissible constant is $D=0$ --- and $2n+0$ is already attained by
part (a). Parts (b) and (c) improve every $n\ge2$, resp.\ $n\ge3$,
pointwise, yet leave the optimal pair $(C,D)=(2,0)$ unchanged: a linear
envelope cannot register these gains. The finer object is the exact envelope
$T(\cdot)$ itself, ordered pointwise.

\emph{2. Which $n$ are hard is a $2$-adic question.}
There is a second exactly fair method, of a different design, with
\[
T(n)\;\le\;n+2^{\nu_2(n)}\qquad\text{for every } n ,
\]
where $2^{\nu_2(n)}$ is the largest power of two dividing $n$: stop at the
first \emph{aligned dyadic interval} $[2Lk+1,2L(k+1)]$ ($L$ a power of two)
whose two halves are constant and distinct, and declare \head{} iff the left
half is the $0$-block. Exchanging the two constant halves of the stopping
interval is again a probability-preserving, verdict-reversing involution on
stopped prefixes (no earlier stopping interval is created or destroyed), so
the method is exactly fair; and if $L=2^{\nu_2(n)}$, the blocks
$[n-L+1,n]$ and $[n+1,n+L]$ are the two halves of an aligned interval, the
left one inside the initial run, so a stop occurs by flip $n+L$. Thus this
method reaches the floor $n+1$ at \emph{every odd} $n$, stays within
$\tfrac43n$ whenever $n$ is not a power of two, but pays $2n$ at
$n=2^k$. Theorem~\ref{thm:c}'s method has the mirror-image profile: fast
($\tfrac32n$) \emph{at} powers of two, slowest ($2n-2$) just above them. The
two are Pareto-incomparable, and every method we know pays slope $2$
somewhere in the vicinity of the powers of two: the linear-envelope
difficulty of the problem concentrates $2$-adically.
\begin{center}
\begin{tabular}{lcccccccccc}
\toprule
$n$ & 1 & 2 & 3 & 4 & 5 & 6 & 7 & 8 & 9 & 10\\
\midrule
floor $n+1$ & 2 & 3 & 4 & 5 & 6 & 7 & 8 & 9 & 10 & 11\\
Theorem~\ref{thm:c} & 2 & 3 & 4 & 6 & 8 & 8 & 8 & 12 & 16 & 16\\
dyadic contrast & 2 & 4 & 4 & 8 & 6 & 8 & 8 & 16 & 10 & 12\\
best known & 2 & 3 & 4 & 6 & 6 & 8 & 8 & 12 & 10 & 12\\
optimal? & yes & yes & yes & ? & yes & ? & yes & ? & yes & ?\\
\bottomrule
\end{tabular}
\end{center}
(``Best known'' minimizes over the two methods separately; no single method
attaining the whole row is known. Odd $n$ are settled at the floor; the
first open value is $T(4)\in\{5,6\}$.)

\emph{3. The frontier.}
Let $C^{*}=\inf\{C:\ T(n)\le Cn+D \text{ for some method and constant }
D\}$. Lemma~\ref{lem:floor} gives $C^{*}\ge1$ and part (a) gives
$C^{*}\le2$; we do not know whether any method beats slope $2$, nor even
whether $C=1$ (that is, $T(n)\le n+D$) is impossible, though we expect it
is. The question has a concrete equivalent form. A method with deadline
$Cn+D$ never stops before the break (Lemma~\ref{lem:floor}), and after the
break at $b^m\bar b$ its stopping words form a complete prefix code of
depth at most $d_m:=(C-1)m+D-1$; padding to exact depth $d_m$ as in
Lemma~\ref{lem:exact} and collecting \head{} words, one checks that such a
method exists if and only if there are polynomials
\[
W_m(p)=\sum_{k=0}^{d_m}w_{m,k}\,p^{\,d_m-k}q^{\,k},\qquad
V_m(p)=\sum_{k=0}^{d_m}v_{m,k}\,p^{\,d_m-k}q^{\,k},
\]
with integer coefficients $0\le w_{m,k},v_{m,k}\le\binom{d_m}k$, satisfying
\[
\sum_{m\ge1}p^{m}q\,W_m(p)\;+\;\sum_{m\ge1}q^{m}p\,V_m(p)\;=\;\frac12
\qquad\text{for all } p\in(0,1).
\]
Part (a) solves this with $d_m=m-1$; whether it is solvable with
$d_m=\gamma m+O(1)$ for some $\gamma<1$ is exactly the proposer's question,
with $C^{*}=1+\inf\gamma$. For $C=1$ the $W_m$ range over a fixed finite
set of polynomials, which is why we expect a rigidity argument to refute
that case; and small instances --- e.g.\ whether $T(4)=5$ is attainable at
all --- are finite feasibility questions. We would be glad to learn what
value of $C$ the proposer has in mind, and whether the intended reading of
$Cn+D$ is uniform or asymptotic.

\end{document}
